

\chapter{Finiteness theorems in \texorpdfstring{$\ell$}{l}-adic cohomology}\label{VII}











\section{Statement of the theorems}\label{VII:1}

Throughout this exposé, $S$ will be a Noetherian scheme and $A$ a left
Noetherian torsion ring annihilated by an integer invertible on $S$. Our main
result is the following.






\begin{theorem_}\label{VII:1-1}
Suppose that $S$ is regular of dimension $0$ or $1$. Let $f:X\to Y$ be a
morphism of $S$-schemes of finite type, and let $\sF$ be a constructible sheaf
of left $A$-modules on $X$. Then the sheaves $\eR^i f_\ast \sF$ are
constructible.
\end{theorem_}

The proof will be given in Section \ref{VII:2} and in \ref{VII:3-10}.






\subsection{Remark}\label{VII:1-2}

In characteristic $0$, this result is less general than
\cite[XIX paragraphe 5]{sga4}, which proves the conclusion of the theorem for
every morphism of finite type of excellent schemes of characteristic $0$.
The proof of \cite[XIX]{sga4} uses, on the one hand, resolution of
singularities and, on the other, the assumption that the schemes have equal
characteristic (in order to deduce the ``purity theorem'' from resolution).






\subsection{Remark}\label{VII:1-3}

We shall say that a complex $K\in\ob\eD(X,A)$ is \emph{constructible} if its
cohomology sheaves are constructible, and shall use a subscript $c$ for the
subcategory of $\eD(X,A)$ (or $\eD^+$, $\eD^-$, $\eD^b$) consisting of
constructible complexes. In the language of derived categories, which we
shall use freely, \ref{VII:1-1} says that
$\eR f_\ast:\eD^+(X,A) \to \eD^+(Y,A)$ sends $\eD_c^+$ into $\eD_c^+$.
More explicitly:
\begin{enumerate}[\indent a)]
  \item For $K$ concentrated in a sheaf $\sF$ in degree $0$, one has
    $\sH^i\eR f_\ast K = \eR^i f_\ast \sF$; the derived statement therefore
    implies the theorem.
  \item Conversely, one invokes the spectral sequence
    \[
      E_1^{p q} = \eR^q f_\ast \sH^p(K) \Rightarrow \sH^{p+q} \eR f_\ast K \text{.}
    \]
\end{enumerate}

Speaking the language of derived categories has the advantage of replacing,
by the simple transitivity formula $\eR(f g)_\ast = \eR f_\ast \eR g_\ast$,
the Leray spectral sequence
$E_2^{p q} = \eR^p f_\ast \eR^q g_\ast \Rightarrow \eR^{p+q} (f g)_\ast$.

Under the hypotheses of \ref{VII:1-1}, it follows from \cite{sga4} that
$\eR f_\ast$ has finite cohomological dimension. This makes it possible above
to replace $\eD^+$ by $\eD$, $\eD^-$, or $\eD^b$.






\subsection{Ideas of the proof}\label{VII:1-4}
\begin{enumerate}[\indent a)]
  \item Poincaré duality treats the case where $X$ is smooth, $\sF$ is
    locally constant, and $Y=S$. The dimension hypothesis is seen na\"ively
    when calculating the $\rHom$ on $S$.
  \item Factor $f$ as $g j$: $g$ proper and $j$ an open immersion. One has
    $\eR f_\ast = \eR g_\ast \eR j_\ast$, and the finiteness theorem for
    proper morphisms controls $\eR g_\ast$. By dévissage, one reduces to
    assuming that $X$ is smooth, $\sF$ is locally constant, $f$ is an open
    immersion, and $Y$ is proper over $S$.
  \item Naively, an induction on $\dim X$ (with a change of $S$) makes it
    possible to suppose that the theorem is true outside a part of $Y$ finite
    over $S$. Denoting the structure morphisms by $a:X\to S$ and $b:Y\to S$,
    one then shows that if the theorem were false for $\sF$ and
    $\eR f_\ast$, it would also be false for $\sF$ and
    $\eR b_\ast \eR f_\ast = \eR a_\ast$: a contradiction.
\end{enumerate}






\begin{corollary_}\label{VII:1-5}
Under the hypotheses of the theorem, the categories $\eD_c(X,A)$ and
$\eD_c(Y,A)$ are taken into one another by the $4$ operations
$\eR f_\ast$, $\eR f_!$, $f^\ast$, $\eR f^!$.
\end{corollary_}

The case of $\eR f_\ast$ is treated above, that of $\eR f_!$ in
\cite[XVII 5.3]{sga4}, and that of $f^\ast$ is clear. It remains to consider
$\eR f^!$. The problem is local, which makes it possible to suppose that $f$
factors as $X\xrightarrow i Z \xrightarrow g Y$ ($i$ a closed immersion and
$g$ smooth, purely of relative dimension $n$). Poincaré duality
$\eR g^! K \iso g^\ast K(n)[2 n]$ and the transitivity
$\eR f^! = \eR i^! \eR g^!$ reduce us to proving \ref{VII:1-5} for $i$.
For every $L\in \eD(Z,A)$, if $j$ is the inclusion into $Z$ of
$U=Z\setminus X$, then $i_\ast \eR i^! L$ is the mapping cylinder of
$L\to \eR j_\ast j^\ast L$, and \ref{VII:1-5} for $i$ follows from
\ref{VII:1-1} for $j$.






\begin{corollary_}\label{VII:1-6}
Let $X$ be as in the theorem, and suppose that $A$ is commutative. Then, if
$\sF$ and $\sG$ are constructible sheaves of $A$-modules on $X$, the
$\Ext_A^i(\sF,\sG)$ are constructible: $\rHom$ sends
$\eD_c^-(X,A)\times \eD_c^+(X,A)$ into $\eD_c^+(X,A)$.
\end{corollary_}

By dévissage of $\sF$, we reduce to assuming that $\sF$ has the form
$j_!\sF_1$ ($j:Y\hookrightarrow X$ a locally closed immersion and $\sF_1$
locally constant on $Y$). One then has
\[
  \rHom(j_! \sF_1,\sG) = \eR j_\ast \Hom(\sF_1,\eR j^! \sG) \text{:}
\]
by \ref{VII:1-3}, we are reduced to the case where $\sF$ is locally
constant---or even constant, since the problem is local. For $\sF$ locally
constant and constructible, one has
$\Hom(\sF,\sG)_x=\hom(\sF_x,\sG_x)$, and likewise for $\rHom$. For $\sF$
constant, the $\Ext^i$ can therefore be calculated using a projective
resolution of finite type of its constant value, and the $\Ext^i$ are
constructible if $\sG$ is---whence the corollary.






\subsection{Remark}\label{VII:1-7}

Since $\eR f_\ast$ has finite cohomological dimension, it takes complexes of
finite tor-dimension (resp. $\leqslant d$) to complexes of finite
tor-dimension (resp. $\leqslant d$) \cite[XVII 5.2.11]{sga4}.
\cite[XVII 5.2.10]{sga4} and the proof of \ref{VII:1-5} then show that finite
tor-dimension is stable under the $4$ operations $\eR f_\ast$,
$\eR f_!$, $f^\ast$, $\eR f^!$. A variant of the argument of
\ref{VII:1-6} (apply dévissage to $K$ along a partition of $X$ so that the
cohomology sheaves become locally constant, then localize so as to replace $K$
by a finite complex of locally free sheaves of finite type) then shows that,
for $A$ commutative, $\rHom$ induces
\[
  \rHom:\eD_\text{ctf}^b(X,A)\times \eD_\text{tf}^b(X,A) \to \eD_\text{tf}^+(X,A)
\]
(t.f.: finite tor-dimension).






\subsection{}\label{VII:1-8}

Under the hypotheses of the theorem, the same method allows us to prove a
local biduality theorem $K\iso D D K$ (Section \ref{VII:4}). We also prove a
finiteness theorem for vanishing-cycle sheaves. For arbitrary $S$, one still
obtains a ``generic'' theorem:






\begin{theorem_}\label{VII:1-9}
Let $f:X\to Y$ be a morphism of $S$-schemes of finite type, and let $\sF$ be
a constructible sheaf of $A$-modules on $X$. There is a dense open subset
$U$ of $S$ such that
\begin{enumerate}[\indent (i)]
  \item Above $U$, the $\eR^i f_\ast \sF$ are constructible, and all but
    finitely many of them vanish.
  \item Formation of the $\eR^i f_\ast \sF$ is compatible with every base
    change $S' \to U\subset S$.
\end{enumerate}
\end{theorem_}

For $S$ the spectrum of a field, one has $U=S$.

If $S$ is the spectrum of a field, a passage-to-the-limit argument gives
compatibility with $S$-base changes for the $\eR^i f_\ast$ for every
quasi-compact quasi-separated morphism of $S$-schemes and every sheaf $\sF$.






\begin{corollary_}\label{VII:1-10}
Let $X$ be a scheme of finite type over a separably closed $k$, and let
$\sF$ be a constructible sheaf of $A$-modules. Then the $\h^i(X,\sF)$ are of
finite type.
\end{corollary_}

This is the special case $S=\spec(k)$, $Y=S$.

\begin{corollary_}\label{VII:1-11}
Let $X$ and $Y$ be two schemes of finite type over a separably closed $k$,
and let $K\in \ob\eD^-(X,A^\circ)$,
$L\in \ob\eD^-(Y,A)$ be complexes of sheaves of right and left modules. The
K\"unneth arrow
\[
  \eR \Gamma(X,K) \lotimes \eR \Gamma(Y,L) \to \eR\Gamma(X\times Y,\pr_1^\ast K \lotimes \pr_2^\ast L)
\]
is an isomorphism.
\end{corollary_}

Proceed as in \cite[XVII 5.4.3]{sga4}:
\[\xymatrix{
  X\times Y \ar[r]^-{\text{pr}_1} \ar[d]^-{\text{pr}_2}
    & X \ar[d]^-a \\
  Y \ar[r]^-b
    & \spec(k)
}\]
base change by $b$ shows that $\eR\pr_{2,\ast}{} \pr_1^\ast K$ is
$b^\ast \eR\Gamma(X,K)$. One then has (cf.
\cite[XVII 5.2.11]{sga4} and the proof)

\begin{align*}
  \eR\Gamma(X\times Y,\pr_1^\ast K\lotimes \pr_2^\ast L)
    &= \eR\Gamma(Y,\eR \pr_{2,\ast}(\pr_1^\ast K\lotimes \pr_2^\ast L)) \\
    &= \eR\Gamma(Y,(\eR\pr_{2,\ast}\pr_1^\ast K)\lotimes L) \\
    &= \eR\Gamma(Y,b^\ast \eR\Gamma(X,K)\lotimes L) \\
    &= \eR\Gamma(X,K)\lotimes \eR\Gamma(Y,L)
\end{align*}






\subsection{}\label{VII:1-12}

Finally, as a local counterpart of \ref{VII:1-9}(ii), we shall prove a
generic local acyclicity theorem (\ref{VII:2-13}).











\section{Generic theorems}\label{VII:2}

In this section, we prove \ref{VII:1-9} and the generic local acyclicity theorem
\ref{VII:2-13}. To prove \ref{VII:1-9}, we immediately reduce to the case where
$S$ is integral. Let $\eta$ be its generic point.





\subsection{}\label{VII:2-1}

We begin by proving \ref{VII:1-9} under the following additional assumptions:
$X$ is smooth over $S$, purely of relative dimension $n$, $A=\dZ/m$, $\sF$ is
locally constant, and $Y=S$. Let $\sF'=\Hom(\sF,\dZ/m)$. We shall show that if
the $\eR^i f_! \sF'$ are locally constant, the conclusions of \ref{VII:1-9}
hold with $U=S$. Recall that if $\sL$ is a constructible locally constant
sheaf, the $\Ext^i(\sL,\sG)$ may be computed fiber by fiber, and are
constructible if $\sG$ is. Recall also that $\dZ/m$ is an injective
$\dZ/m$-module and is the dualizing module, so that
$\sF=\rHom(\sF',\dZ/m)$. Poincaré duality
\[
  \eR f_\ast \rHom(K,\eR f^! L) = \rHom(\eR f_! K,L) \text{,} 
\]
for $K=\sF'$, $L=\dZ/m$, $\eR f^! L = \dZ/m[2n](n)$, therefore gives
\[
  \eR^{2n-i} f_\ast \sF = \Hom(\eR^i f_! \sF',\dZ/m)(-n) \text{.} 
\]

The sheaf on the right is locally constant and its formation commutes with
every base change, which proves the assertion.





\subsection{}\label{VII:2-2}

Let us prove \ref{VII:1-9} under the additional assumptions that $X$ is smooth
over $S$, $\sF$ is locally constant, and $Y=S$.
\begin{enumerate}[a)]
  \item By decomposing $X$ into connected components, we reduce to the case
    where it is purely of some relative dimension $n$ over $S$.
  \item By decomposing $A$ as a product, we reduce to the case where
    $\ell^m A=0$, with $\ell$ a prime invertible on $S$; we then filter $\sF$
    by the $\ell^k \sF$ and use the corresponding spectral sequence to reduce
    to the case $\ell\sF=0$. We may then replace $A$ by $A/\ell A$ and assume
    that $\ell A=0$.
  \item After replacing $S$ by a suitable $U$, we may assume that there exists
    a finite étale Galois covering $X_1/X$, with Galois group $G$, such that
    the inverse image $\sF_1$ of $\sF$ on $X_1$ is a constant sheaf, with
    constant value $F$. Writing $f_1$ for the projection of $X_1$ onto $S$,
    we then have
    \[
      \eR^i {f_1}_\ast \sF_1 = (\eR^i {f_1}_\ast \dZ/\ell)\otimes_{\dZ/\ell} F \text{.} 
    \]
    We also have the Hochschild--Serre spectral sequence (or the Leray
    spectral sequence for the covering $X_1/X$)
    \[
      \underline\h^p(G,\eR^q {f_1}_\ast \sF_1) \Rightarrow \eR^{p+q} f_\ast \sF \text{.} 
    \]
    By \ref{VII:2-1}, after shrinking $U$ if necessary, we may assume that
    the $\eR^i {f_1}_\ast \dZ/\ell$ are locally constant and that their
    formation commutes with every base change. The same property then holds
    for the $\eR^i f_\ast \sF$. Finally, if $\bar\eta$ is a geometric point
    localized at the generic point $\eta$ of $S$, almost all the
    $\eR^i f_\ast\sF$ vanish, because the
    $(\eR^i f_\ast \sF)_{\bar\eta} = \h^i(X_{\bar \eta},\sF)$ do.
\end{enumerate}





\subsection{}\label{VII:2-3}

We prove by induction on $n$ that
\begin{equation*}\tag{$\ast_n$}\label{VII:eq:2-3-1}
  \begin{array}{c}
    \text{The conclusions of \ref{VII:1-9} hold when $\dim{X_\eta}\leqslant n$} \\
    \text{and $f$ is an open immersion with dense image.}
  \end{array}
\end{equation*}

For $n=0$, after shrinking $S$ if necessary, we have $X=Y$: $(\ast_0)$ is
clear. Assume $(\ast_{n-1})$, and prove \eqref{VII:eq:2-3-1}. In
($\ast_{n-1}$), ``open immersion'' may be replaced by ``immersion,'' as is
seen by factoring it as an open immersion followed by a closed immersion.





\begin{lemma_}\label{VII:2-4}
After shrinking $S$ if necessary, the conclusions of \ref{VII:1-9} hold over
$Y'\subset Y$, where the complement $Y_1$ of $Y'$ is finite over $S$.
\end{lemma_}

The assertion is local on $Y$, which we may assume affine:
$Y\subset \dA_S^n$. The induction hypothesis ($\ast_{n-1}$) applies to
\[\xymatrix{
  X \ar[r]^-f \ar[dr] 
    & Y \ar[d]^-{\pr_i} \\
  & \dA_S^1 
}\]
Thus, for each $i$, there exists a dense open subset $U_i$ of $\dA_S^1$ such
that the conclusions of \ref{VII:1-9} hold over $\pr_i^{-1}(U_i)$; they hold
over the union of the $\pr_i^{-1}(U_i)$, and \ref{VII:2-4} follows.





\addtocounter{subsection}{1} % a section is skipped in the original
\subsection{}\label{VII:2-6}

Let us prove \eqref{VII:eq:2-3-1} when $X$ is smooth over $S$ and $\sF$ is
locally constant on $X$. Since the problem is local on $Y$, we may first
assume $Y$ affine and then projective (replacing $Y$ by its closure in a
projective space). Let $i:Y_1\hookrightarrow Y$ and $j:Y' \to Y$ be as in
\ref{VII:2-4}.
\[\xymatrix{
  X \ar@{^{(}->}[r] \ar[dr]_-a 
    & Y \ar[d]^-b 
    & \ar[l]_-i \ar[dl]^-{b_1} Y_1 \\
  & S 
}\] 
After shrinking $S$ if necessary, we know that $j^\ast \eR f_\ast \sF$ is
constructible and that its formation commutes with every base change in $S$;
we also know that $\eR a_\ast \sF = \eR b_\ast \eR f_\ast \sF$ is
constructible and that its formation commutes with every base change.

Apply $\eR b_\ast$ to the triangle defined by the exact sequence
\begin{equation*}\tag{1}\label{VII:eq:2-6-1}
\xymatrix{
  0 \ar[r] 
    & j_! j^\ast \eR f_\ast \sF \ar[r] 
    & \eR f_\ast \sF \ar[r] 
    & i_! i^\ast \eR f_\ast \sF \ar[r] 
    & 0 \text{:} 
}
\end{equation*}
we obtain a triangle
\begin{equation*}\tag{2}\label{VII:eq:2-6-2}
\xymatrix{
  \ar[r] 
    & \eR b_\ast j_! j^\ast \eR f_\ast \sF \ar[r] 
    & \eR a_\ast \sF \ar[r] 
    & {b_1}_\ast i^\ast \eR f_\ast \sF \ar[r] 
    & 
}
\end{equation*}
in which the first two terms are constructible and their formation commutes
with every base change in $S$ (for the first, by the finiteness theorem for
the proper morphism $b$). The same therefore holds for the third. Since $b_1$
is finite, it follows that $i^\ast \eR f_\ast \sF$ is constructible and its
formation commutes with every base change in $S$; the same holds for
$\eR f_\ast \sF$ by \eqref{VII:eq:2-6-1}.





\subsection{}\label{VII:2-7}

Let us prove \eqref{VII:eq:2-3-1} in general. We first reduce to the case where
$X$ contains a dense open subset $V$ smooth over $S$. When $S$ is the spectrum
of a perfect field, it is enough to replace $X$ by $X_\text{red}$ (and $Y$ by
$Y_\text{red}$). In general, we must shrink $S$, make a finite, radicial,
surjective base change $S' \to S$, and replace $X$ and $Y$ by $X_\text{red}$
and $Y_\text{red}$. Since the étale topology is insensitive to finite,
radicial, surjective morphisms, this is harmless. After shrinking $V$, we may
assume that $\sF$ is locally constant on $V$.
\[\xymatrix{
  V \ar@{^{(}->}[r]^j 
    & X \ar@{^{(}->}[r]^-f 
    & Y \text{.} 
}\]
Define $\Delta$ by the triangle
\begin{equation*}\tag{1}\label{VII:eq:2-7-1}
\xymatrix{
  \ar[r] 
    & \sF \ar[r] 
    & \eR j_\ast j^\ast \sF \ar[r] 
    & \Delta \ar[r] 
    & \text{.} 
}
\end{equation*}
The cohomology sheaves of $\Delta$ are supported on $X\setminus V$, and
$\dim(X\setminus V)_\eta < n$. The induction hypothesis therefore allows us
to assume that $\eR f_\ast \Delta$ is constructible. Applying $\eR f_\ast$ to
the triangle \eqref{VII:eq:2-7-1}, we obtain a triangle
\[\xymatrix{
  \ar[r] 
    & \eR f_\ast \sF \ar[r] 
    & \eR(f j)_\ast j^\ast \sF \ar[r] 
    & \eR f_\ast \Delta \ar[r] 
    & 
}\]
in which two of the vertices are constructible and their formation commutes
with every base change in $S$. The same therefore holds for the third.





\subsection{}\label{VII:2-8}

Let us prove \ref{VII:1-9}. The problem is local on $Y$, which we may assume
affine. Taking an affine covering of $X$ and invoking the Leray spectral
sequence, we reduce to the case where $X$ is affine as well. This ensures that
we can factor $f$ as an open immersion followed by a proper morphism:
$f=g j$, hence $\eR f_\ast = \eR g_\ast \eR j_\ast$. The open immersion is
covered by \eqref{VII:eq:2-3-1}, and the proper morphism by the finiteness
theorem.

The following corollaries are proved as in \hyperref[VII:1]{Section 1}.





\begin{corollary_}\label{VII:2-9}
Under the hypotheses of the theorem, for $K$ in $\eD_c^b(X,A)$ or
$\eD_c^b(Y,A)$ respectively, there exists a nonempty open subset $U$ of $S$
over which $\eR f_\ast K$, $\eR f_! K$, $f^\ast K$, and $\eR f^! K$ belong to
$\eD_c^b(Y,A)$ or $\eD_c^b(X,A)$, and their formation commutes with every base
change $S' \to U\subset S$.
\end{corollary_}





\begin{corollary_}\label{VII:2-10}
Let $X$ be as in the theorem, and suppose that $A$ is commutative. Then, if
$\sF$ and $\sG$ are constructible sheaves of $A$-modules on $X$, over a dense
open subset $U$ of $S$, the $\Ext_A^i(\sF,\sG)$ are again constructible and
their formation commutes with every base change $S' \to U\subset S$.
\end{corollary_}





\subsection{}\label{VII:2-11}

If $x$ is a geometric point of a scheme $X$, we denote by $X_x$ the strict
henselization of $X$ at $x$. For $f:X\to S$ and $t$ a geometric point of $S$,
we denote by $X_t$ the geometric fiber of $X$ at $t$. Finally, for $x$ a
geometric point of $X$ and $t$ a geometric point of $S_{f(x)}$, $(X_x)_t$ is
the fiber at $t$ of $X_x \to S_{f(x)}$.





\begin{definition_}\label{VII:2-12}
Let $f:X\to S$ and $K\in \ob\eD^+(X,A)$. We say that $f$ is \emph{locally
acyclic at $x$ relative to $K$} if, for every geometric point $t$ of
$S_{f(x)}$, we have $K_x=\eR\Gamma(X_x,K)\iso \eR\Gamma(X_{x,t},K)$. We say
that $f$ is \emph{locally acyclic relative to $K$} if this holds for every
$x$, and \emph{universally locally acyclic relative to $K$} if it remains true
after every base change $S' \to S$.
\end{definition_}





\begin{theorem_}\label{VII:2-13}
Let $f:X\to S$ be a morphism of finite type and let $\sF$ be constructible on
$X$. There exists a dense open subset $U$ of $S$ over which $f$ is universally
locally acyclic relative to $\sF$.
\end{theorem_}

We shall use the following result, proved in the appendix.





\begin{lemma_}\label{VII:2-14}
Let $X\xrightarrow f S_1 \xrightarrow g S_2$. If $g$ is smooth and $f$ is
universally locally acyclic relative to $K$, then so is $g f$.
\end{lemma_}

We reduce to the case where $S$ is integral, with generic point $\eta$, and
argue by induction on $\dim{X_\eta}$. We first deduce from the induction
hypothesis that
\begin{equation*}\tag{A}\label{VII:eq:2-14-1}
  \begin{array}{c}
    \text{After shrinking $S$, there exists $T\subset X$, finite over $S$, such that $f$} \\
    \text{is universally locally acyclic outside $T$.} 
  \end{array}
\end{equation*}

Since this question is local, we localize on $X$ and factor $f$ as
$X\xrightarrow u \dA_S^1 \to S$. The induction hypothesis applies
to $u$, and we conclude by \ref{VII:2-14} and the usual arguments.

To prove the theorem, we may assume that $X$ is proper, since the problem is
local. We shrink $S$ so that the $\eR^i f_\ast \sF$ are locally constant and
\eqref{VII:eq:2-14-1} applies, and use the following lemma.





\begin{lemma_}\label{VII:2-15}
Let $f:X\to S$ be proper and let $\sF$ be constructible. Suppose that the
$\eR^i f_\ast \sF$ are locally constant and that, outside $T\subset X$ finite
over $S$, $f$ is locally acyclic relative to $\sF$. Then $f$ is locally
acyclic relative to $\sF$.
\end{lemma_}

Let $s$ be a geometric point of $S$, let $S_s$ be the strict localization of
$S$ at $s$, and let $t$ be a geometric point of $S_s$. Let $X_{(s)}$ be the
inverse image of $X$ on $S_s$, let $\bar j:X_t \to X_{(s)}$, and let
$i:X_s \hookrightarrow X_{(s)}$. Again denoting by $\sF$ the inverse image of
$\sF$ on $X_{(s)}$ or $X_t$, we must prove that
$i^\ast \sF \iso i^\ast \eR \bar j_\ast \sF$. Let $\Delta$ be the mapping
cylinder of this map. Its cohomology sheaves are supported on $T_s$. To prove
that $\Delta=0$, it therefore suffices to prove that
$\eR\Gamma(X_s,\Delta)=0$. This is the mapping cylinder of
$(\eR f_\ast \sF)_s = \eR\Gamma(X_s,\sF) \to \eR\Gamma(X_s,i^\ast \eR\bar j_\ast \sF) = \eR\Gamma(X_{(s)},\eR\bar j_\ast \sF) = \eR\Gamma(X_t,\sF) = (\eR f_\ast \sF)_t$.
By hypothesis this specialization morphism is an isomorphism, and
$\Delta=0$.





\begin{corollary_}\label{VII:2-16}
When $S$ is the spectrum of a field, every $S$-scheme $X$ is universally
locally acyclic relative to $K$, for every $K\in\ob\eD^+(X,A)$.
\end{corollary_}

This follows from \ref{VII:2-13} by passage to the limit (which is possible
because the $U$ of \ref{VII:2-13} is always equal to $S$.








\section{Proof of \texorpdfstring{\ref{VII:1-1}}{1.1} and constructibility of the vanishing-cycle sheaves}\label{VII:3}






\subsection{}\label{VII:3-1}

Let $S$ be a strictly local trait (the spectrum of a henselian discrete
valuation ring with separably closed residue field). Denote its closed and
generic points by $s$ and $\eta$, and let $\bar\eta$ be a geometric point
localized at $\eta$ (a separable closure of $k(\eta)$). Let $X$ be over $S$
and let $\sF$ be a sheaf on $X_\eta$. Recall the definition of the
vanishing-cycle sheaves $\eR^i \Psi_\eta(\sF)$ (sheaves on $X_s$, endowed with
an action of $\gal(\bar\eta/\eta)$). Let $i:X_s\hookrightarrow X$,
$j:X_\eta \hookrightarrow X$, and let $\bar j$ be the composite
$X_{\bar\eta}\to X_\eta \hookrightarrow X$; then
\[
  \eR^i \Psi_\eta(\sF) = i^\ast \eR^i \bar j_\ast (\bar j^\ast \sF) \text{.} 
\]

Here is the main theorem of this subsection. It improves
\cite[XIII 2.3.1, 2.4.2]{sga7}.






\begin{theorem_}\label{VII:3-2}
Suppose that $X$ is of finite type over $S$ and that $\sF$ is constructible.
Then the vanishing-cycle sheaves $\eR^i \Psi_\eta(\sF)$ are constructible.
\end{theorem_}

We argue by induction on the dimension $n$ of $X_\eta$. We may moreover
assume---and do assume---that $X_\eta$ is dense in $X$ (replacing $X$ by
$\bar X_\eta$ does not change the $\eR^i \Psi$, which are supported on
$\bar X_\eta$).






\begin{lemma_}\label{VII:3-3}
If \ref{VII:3-2} holds when $X_\eta<n$, and if
$\dim X_\eta=n$, then there exist constructible subsheaves $\sG_i$ of the
$\eR^i \Psi_\eta(\sF)$ such that the supports of the local sections of
$\eR^i\Psi_\eta(\sF)/\sG_i$ are finite.
\end{lemma_}

Let $s'$ be a geometric generic point of the affine line over $s$, and let
$S'$ be the strict localization at $s'$ of the affine line $\dA_S^1$ over
$S$. It is again a strictly local trait, and uniformizers for $S$ are
uniformizers on $S'$.
\[\xymatrix{
  & \spec(k') \ar[rr] \ar[dd] \ar[dr] 
    & & (S',\eta',s') \ar[dr] \ar[dd] \\
  \bar\eta' \ar[ur] \ar[dr] 
    & & \dA_{\bar S}^1 \ar[d]  
    & & \dA_S^1 \ar[dl] \\
  & \bar\eta\ar@{^{(}->}[r] 
    & \bar S \ar[r] 
    & (S,\eta,s) 
}\]
Let $k'=k(\bar\eta)\otimes_{k(\eta)}k(\eta')$. If $\bar S$ is the
normalization of $S$ in $\bar\eta$, then $k'$ is the fraction field of the
strict localization of $\dA_{\bar S}^1$ at $s'$; it is therefore a field, and
$\gal(\bar\eta/\eta)=\gal(k'/\eta')$. Let $\bar\eta'$ be the spectrum of an
algebraic closure of $k'$. The Galois group $P=\gal(\bar\eta'/k')$ is a
pro-$p$ group, where $p$ is the characteristic exponent of $k(s)$.






\begin{lemma_}\label{VII:3-4}
For every scheme $X'$ over $S'$, and every sheaf $\sF$ on
$X_\eta'=X_{\eta'}'$, the vanishing-cycle sheaves for $X'/S'$ and for $X'/S$
are related by
$\eR^i \Psi_\eta(\sF) =\eR^i\Psi_{\eta'}(\sF)^P$.
\end{lemma_}

The following diagram compares the morphisms used in the definition of
$\eR\Psi$ for $X'/S'$ and $X'/S$:
\[\xymatrix{
  X_{\bar\eta'} \ar[rr] \ar[d] 
    & & X_{\eta'}' \ar@{=}[dd] \ar@{^{(}->}[dr] 
    & & S' \\ 
  X_{\eta'}\times_{\eta'}\spec(k') \ar@{=}[d] \ar[urr] 
    & & & X \ar[ur] \ar[dr] \\
  X_\eta' \ar[rr] 
    & & X_\eta' \ar@{^{(}->}[ur] 
    & & S
}\]
and the lemma follows from the Hochschild--Serre spectral sequence, since $P$
is a pro-$p$ group and $p$ is invertible in $A$.

Let us prove \ref{VII:3-3}. The question is local; this allows us to assume
that $X$ is affine, $X\subset \dA_S^n$. Let $f$ be one of the projections
$X\subset \dA_S^n\to \dA_S^1$. Let $X'$ be the ``localization''
$X_{\dA_S^1} S'$ of $X$, and let $\sF'$ be the inverse image of $\sF$ on $X'$.
\[\xymatrix{
  \sF'\text{ on }X' \ar[r] \ar[d]^-\lambda 
    & S' \ar[d]^-\lambda \\
  \sF \text{ on }X \ar[r]^-f 
    & \dA_S^1 \ar[r] 
    & S \text{.} 
}\]
On $X_s'=X_{s'}'$, we have
\[
  \lambda^\ast \eR^i\Psi_\eta(\sF) = \eR^i \Psi_\eta(\sF') = \eR^i\Psi_{\eta'}(\sF')^P \text{:} 
\]
this sheaf is constructible because the induction hypothesis applies to
$X'/S'$, and we use the following lemma, applied to
$X_s\subset \dA_S^n$ and to the vanishing-cycle sheaves.






\begin{lemma_}\label{VII:3-5}
Let $\dA^n$ be affine $n$-space over a field $k$, let $X\subset \dA^n$, let
$\sF$ be a sheaf of $A$-modules on $X$, and let $\bar\eta$ be a geometric
generic point of $\dA^1$. Let $X_{\bar\eta,i}$ be the geometric generic fiber
of $X\subset \dA^n\xrightarrow{\pr_i} \dA^1$, and let
$\sF_{\bar\eta,i}$ be the restriction of $\sF$ to $X_{\bar\eta,i}$. If the
$\sF_{\bar\eta,i}$ are constructible, there exists a constructible
$\sF'\subset \sF$ such that the local sections of $\sF/\sF'$ have finite
support.
\end{lemma_}

By passage to the limit, for each $i$
\begin{enumerate}[a)]
  \item There exists an étale neighborhood $U$ of $\bar\eta$ and a
    constructible sheaf $\sH$ on
    $X_{U,i} = X\times_{\dA^1,\pr_i} U$ whose restriction
    $\sF_{\bar\eta,i}$ to $X_{\bar\eta,i}$ is the given one.
  \item For a suitable $U$, the isomorphism
    $\sH_{\bar\eta,i} \iso \sF_{\bar\eta,i}$ comes from
    $u:\sH\to \sF_{U,i}$, where $\sF_{U,i}$ is the inverse image of $\sF$ on
    $X_{U,i}$. Let $\varphi:X_{U,i} \to X$. The morphism $u$ defines
    $v_i:\varphi_! \sH \to \sF$, with image $\sF_i'$. We have
    $(\sF/\sF_i')_{\bar\eta,i} = 0$, and for $\sF'$ we take the sum of the
    $\sF_i'$.
\end{enumerate}






\subsection{}\label{VII:3-6}

Let us prove \ref{VII:3-2} in dimension $n$. We may assume that $X$ is affine
(because the problem is local on $X_s$), and then projective over $S$ (by
taking the closure of $X$ in a projective embedding). We then have a spectral
sequence \cite[I 2.2.3]{sga7} (the Leray spectral sequence for $\bar j$, taking
into account the base-change theorem for the proper morphism $X\to S$).
\begin{equation*}\tag{1}\label{VII:eq:3-6-1}
  E_2^{p q} = \h^p(X_s,\eR^q\Psi_\eta(\sF)) \Rightarrow \h^{p+q}(X_{\bar\eta},\sF) \text{.} 
\end{equation*}
Let $\sG^q\subset \eR^q \Psi_\eta(\sF)$ be as in \ref{VII:3-3}, and let
$\sH^q$ be the quotient sheaf. We have $\h^i(X_s,\sH^q)=0$ for $i\ne 0$, and
for $\sH^q$, hence $\eR^q\Psi_\eta(\sF)$, to be constructible, it suffices
that $\h^0(X_s,\sH^q)$ be of finite type.

Compute \eqref{VII:eq:3-6-1} modulo modules of finite type (i.e. in the
quotient of the category of $A$-modules by the thick subcategory of
$A$-modules of finite type). The long cohomology exact sequence deduced from
$0 \to \sG^q \to \eR^q\Psi_\eta(\sF) \to \sH^q \to 0$, together with
\ref{VII:1-10} applied to $\sG^q$, gives
$E_2^{p q} \sim \h^p(X_s,\sH^q)$. In particular,
$E_2^{p q}\sim 0$ for $p\ne 0$,
$E_2^{0 q}\sim \h^q(X_{\bar\eta},\sF) \sim 0$, and this completes the proof.

Parallel arguments give the following result, which improves
\cite[XIII 2.1.12, 2.4.2]{sga7}.






\begin{proposition_}\label{VII:3-7}
The formation of the vanishing-cycle sheaves commutes with changes of traits.
\end{proposition_}

Let $g:(S',\eta',s') \to (S,\eta,s)$ be a (surjective) morphism of strictly
local traits, let $X/S$, let $\sF$ be on $X_\eta$, of torsion prime to the
residual characteristic, and let $(X',\sF')$ be their inverse images on $S'$.
Again denoting by $g$ the morphisms parallel to $g$, such as
$X_{s'}' \to X_s$, we must prove that
\begin{equation*}\tag{3.7.1}\label{VII:eq:3-7-1}
\xymatrix{
  g^\ast \eR\Psi_\eta(\sF) \ar[r]^-\sim 
    & \eR\Psi_{\eta'}(\sF') \text{.} 
}
\end{equation*}
This morphism is defined by the commutative diagram
\[\xymatrix{
  X_{s'}' \ar@{^{(}->}[r]^-{i'} \ar[d]^-g 
    & X' \ar[d]^-g 
    & X_{\eta'}' \ar[l]_-{\bar j'} \ar[d]^-g \\
  X_s \ar@{^{(}->}[r]^-i 
    & X 
    & \ar[l]_-{\bar j} X_{\bar\eta} 
}\]

A passage to the limit reduces us to the case where $X$ is of finite type
over $S$, and we argue by induction on $\dim{X_\eta}$. Since the question is
local, we may assume that $X$ is affine, and then projective. Assume
\eqref{VII:eq:3-7-1} when $\dim{X_\eta}<n$. When $\dim{X_\eta}=n$, it then
follows from the following two assertions, where $\Delta$ is the mapping
cylinder of \eqref{VII:eq:3-7-1}.

\begin{enumerate}[(A)]
  \item The support of the local sections of the cohomology sheaves of
    $\Delta$ is finite.
    \phantomsection\hypertarget{VII:3-A}{}
  \item $\eR\Gamma(X_s,\Delta)=0$, hence $\Delta=0$, by
    (\hyperlink{VII:3-A}{A}).
    \phantomsection\hypertarget{VII:3-B}{}
\end{enumerate}





\subsection{}\label{VII:3-8}

The proof of (\hyperlink{VII:3-A}{A}) parallels that of \ref{VII:3-3}. We
localize on $X$ and consider maps $X\to \dA_S^1$; denoting by $S(x)$ the
strict henselization of $\dA_S^1$ at the generic point of the special fiber,
we apply the induction hypothesis to $X_1/S(x)$
($X_1=X\times_{\dA^1} S(x)$) and to $S'(x)$. Taking invariants under a pro-$p$
group, we find that $\Delta$ vanishes on the geometric generic fiber of
$X_{s'} \to \dA_s^1$. Since this holds locally on $X$ and for every
projection, we obtain (\hyperlink{VII:3-A}{A}).






\subsection{}\label{VII:3-9}

For (\hyperlink{VII:3-B}{B}), the properness of $X$ over $S$ ensures that
\[
  \eR\Gamma(X_{s'}',\eR\Psi_{\eta'}(\sF')) = \eR\Gamma(X_{\bar\eta}',\sF') 
\]
and
\[
  \eR\Gamma(X_s,\eR\Psi_\eta(\sF)) = \eR\Gamma(X_{\bar\eta},\sF) \text{.} 
\]
We have
$\eR\Gamma(X_s,\eR\Psi_\eta(\sF))\iso \eR\Gamma(X_{s'}',g^\ast\eR\Psi_\eta(\sF))$
(invariance under change of the separably closed base field); thus
$\eR\Gamma(X_{s'}', \Delta)$ is the mapping cylinder of
$\eR\Gamma(X_{\bar\eta},\sF) \to \eR\Gamma(X_{\bar\eta}',\sF')$, an
isomorphism by the same invariance theorem.






\subsection{}\label{VII:3-10}

Let us prove \ref{VII:1-1}. We reduce to the case where $S$ is connected,
hence integral. If $\dim S=0$ ($S$ is the spectrum of a field), apply
\ref{VII:1-9}, \ref{VII:1-10}. If $S$ has dimension $1$, Theorem
\ref{VII:1-9} ensures that the $\eR^i f_\ast \sF$ are constructible over the
complement of a finite set $T$ of points of $S$. It remains to see that their
restrictions to the $Y_t$ ($t\in T$) are constructible. It suffices to see
this after localization at $t$: we may assume that $S$ is a strictly local
trait. The essential case is the following.






\begin{lemma_}\label{VII:3-11}
Let $X$ be of finite type over $S$, let $j:X_\eta\hookrightarrow X$,
$i:X_s\hookrightarrow X$, and let $\sF$ be constructible on $X_\eta$. Then
$i^\ast\eR j_\ast \sF$ is constructible.
\end{lemma_}

Let $I=\gal(\bar\eta/\eta)$ be the inertia group. It is known to be an
extension of $\widehat\dZ_{p'}(1)$ by a pro-$p$ group $P$ ($p$ being the
characteristic exponent of the residue field). The Hochschild--Serre spectral
sequence gives
\[
  E_2^{p q} = \sH^p(I,\eR^q\Psi_\eta(\sF))\Rightarrow i^\ast \eR^{p+q}j_\ast\sF \text{.} 
\]
Let $R_t^q = \eR^q\Psi_\eta(\sF)^P$ ($t$ for tame). If $\sigma$ is a
generator of $\widehat\dZ_{p'}(1)$, then
$E_2^{0 q} = \ker(\sigma-1,R_t^q)$,
$E_2^{1 q} = \operatorname{coker}(\sigma-1,R_t^q)$, and $E_2^{p q} = 0$ for
$p\ne 0,1$. These sheaves are constructible because the vanishing-cycle
sheaves are, and \ref{VII:3-11} follows.






\subsection{}\label{VII:3-12}

Let us treat the general case, with $S$ still a strictly henselian trait. If
$X=X_\eta$, then $f$ is the composite $X\to Y_\eta \to Y$, and we apply
\ref{VII:1-9} and \ref{VII:3-11}. In the general case, let
$j:X_\eta\hookrightarrow X$ and let $\Delta$ be the mapping cylinder of
$\sF \to \eR j_\ast j^\ast \sF$. Its cohomology sheaves are supported on
$X_s$; by \ref{VII:1-9} applied to $X_s \to Y_s$,
$\eR f_\ast \Delta$ is therefore constructible. The complex
$\eR f_\ast \eR j_\ast j^\ast \sF = \eR(f j)_\ast j^\ast \sF$ is also
constructible, and we conclude from the triangle
\[\xymatrix{
  \ar[r] 
    & \eR f_\ast \sF \ar[r] 
    & \eR f_\ast \eR j_\ast j^\ast \sF \ar[r] 
    & \eR f_\ast \Delta \ar[r] 
    & \text{.} 
}\]










\section{Local biduality}\label{VII:4}






\subsection{}\label{VII:4-1}

Let $S$ be regular of dimension $0$ or $1$, and put $K_S=A$ (the constant
sheaf with value $A$). For $K\in\ob\eD_\text{ctf}^b(S,A)$, put
$D K=\rHom(K,K_S)\in\ob\eD^+(S,A^\circ)$ (a complex of sheaves of right
$A$-modules). We still have $D K\in\ob\eD_\text{ctf}^b(S,A^\circ)$, and an explicit
calculation, possible because $S$ has dimension $1$, shows that
$K\iso D D K$ (for $A=\dZ/n$, \hyperref[V]{Duality} \ref{V:1-4}).





\subsection{}\label{VII:4-2}

We assume that $A$ is commutative, an assumption that is probably
unnecessary. For $a:X\to S$ of finite type over $S$, put
$K_X=\eR a^! K_S$. For $K\in\ob\eD_\text{ctf}^b(X,A)$, put
$D K=\rHom(K,K_X)\in\ob\eD_\text{ctf}^b(X,A)$. Our main result is the
following.





\begin{theorem_}\label{VII:4-3}
We have $K\iso D D K$ ($K\in \ob\eD_\textnormal{ctf}^b(X,A)$).
\end{theorem_}





\begin{lemma_}\label{VII:4-4}
If $a$ is proper, then $\eR a_\ast K \iso \eR a_\ast D D K$.
\end{lemma_}

Poincaré duality
\[
  \eR a_! \rHom(K,\eR a^! K_S) = \rHom(\eR a_! K,K_S) 
\]
gives us an isomorphism $\eR a_\ast D = D \eR a_! = D \eR a_\ast$. The lemma
follows from the commutativity of the diagram (\hyperref[V]{Duality},
\ref{V:1-2})
\[\xymatrix{
  \eR a_\ast K \ar[r] \ar[d]^-\wr 
    & \eR a_\ast D D K \ar[d]^-\wr \\
  D D \eR a_\ast K \ar[r]^-\sim 
    & D \eR a_\ast D K \text{.} 
}\]





\subsection{}\label{VII:4-5}

By localization, it is enough to treat the cases where $S$ is the spectrum
of a field or a trait. When $S$ is the spectrum of a field, we argue by
induction on $\dim X$. Since the problem is local, we may assume that $X$ is
proper. Moreover, the induction hypothesis, by the usual arguments, ensures
that the mapping cylinder $\Delta$ of $K\to D D K$ has skyscraper cohomology
sheaves. The lemma then ensures that $\eR a_\ast \Delta=0$, hence that
$\Delta=0$.





\subsection{}\label{VII:4-6}

When $S$ is a trait, we already know that \ref{VII:4-3} holds on the generic
fiber (\ref{VII:4-5}). We may also assume that $X$ is proper over $S$, and we
argue by induction on $\dim{X_s}$, with $X$ proper over $S$. The mapping
cylinder $\Delta$ again has skyscraper cohomology sheaves on $X_s$, and we
conclude as above.





\subsection{}\label{VII:4-7}

Suppose that $A=\dZ/m$. We then have a variant of the preceding theory in
which $K$ is arbitrary in $\eD_c^b(X,A)$. The complex $\eR a^! \dZ/m$ has
finite injective dimension \cite[XVIII 3.1.7]{sga4}. This ensures that $D K$
still belongs to $\eD_c^b(X,A)$. To prove local biduality over $S$, we use the
fact that $\dZ/m$ is the dualizing $\dZ/m$-module. We then proceed as above.










\section{Appendix, by L.\ Illusie}\label{VII:5}

In this appendix, which takes up certain parts of the late
\cite[II]{sga5}, and a letter from P.\ Deligne to the author, we prove
\ref{VII:2-14}, as well as various generalizations and variants of the
specialization theorems \cite[XVI 2.1]{sga4} and
(\hyperref[I]{Étale cohomology} \ref{I:5-1-7}).

The schemes and morphisms considered will be assumed quasi-compact and
quasi-separated. If $X$ is a scheme and $x$ a geometric point of $X$, we
denote by $X(x)$ the strict localization of $X$ at $x$.

Fix a base scheme $S$ and a sheaf of rings $A$ on $S$ (not necessarily
commutative or noetherian). If $X$ is an $S$-scheme, we write $\eD(X)$ for
$\eD(X,A_X)$.





\subsection{Cohomological properness}\label{VII:5-1}






\begin{proposition}\label{VII:5-1-1}
Let $f:X\to Y$ be an $S$-morphism and let $E\in\ob \eD^+(X)$. The following
conditions are equivalent:
\begin{enumerate}[\indent (i)]
  \item The formation of $\eR f_\ast E$ commutes with every finite base change
    $S'\to S$.
  \item The formation of $\eR f_\ast E$ commutes with every quasi-finite base
    change (or inverse limit of quasi-finite morphisms) $S'\to S$.
  \item For every specialization arrow $i:t\to S(s)$
    (\hyperref[I]{Étale cohomology}, \ref{I:5-1-2}), if $\bar S$ denotes
    the normalization in $k(\bar t)$ of the integral scheme that is the
    closure of $i(t)$ in $S(s)$ (by
    \cite[IV, 6.15.5, 18.6.11, 18.8.16]{ega}, $\bar S$ is therefore local and
    integral, with closed point radicial over $s$), if
    $\bar f:\bar X\to \bar Y$ is the morphism deduced from $f$ by the base
    change $\bar S\to S$, and if $f_s:X_s\to Y_s$ is the fiber of $f$ at $s$,
    then the base-change arrow
    \[
      (\eR \bar f_\ast (E|\bar X))|Y_s \to \eR {f_s}_\ast(E|X_s) 
    \]
    is an isomorphism.
\end{enumerate}
\end{proposition}

The equivalence of (i) and (ii) follows from Zariski's Main Theorem, and it is
clear that (ii) implies (iii). Let us prove that (iii) implies (ii). For every
$S'\to S$ that is an inverse limit of quasi-finite morphisms, and every
geometric point $s$ of $S'$, denote by $C(S',s)$ the cone of the base-change
arrow
\[
  (\eR f'_\ast(E|X'))|Y_s \to \eR{f_s}_\ast (E|X_s) \text{,} 
\]
where $f':X'\to Y'$ is deduced from $f$ by the base change $S'\to S$, and
$f_s:X_s\to Y_s$ is the fiber of $f$ at $s$. We must prove that, for every
$(S',s)$ as above, $C(S',s)$ is acyclic. We shall prove by induction on $N$ that
$\h^i C(S',s) = 0$ for $i\leqslant N$. By the Main Theorem, passage to the
limit, and elimination of radicial morphisms, we may assume that $S$ and $S'$
are strictly local, that $S'$ is integral with closed point $s$, and that
$S'\to S$ is finite. Let $t$ be a geometric generic point of $S'$, and let
$\bar S$ be the normalization of $S'$ in $k(t)$. Let $\bar S_\bullet$ be the
simplicial scheme that is the coskeleton of $\bar S\to S'$
($\bar S_0=\bar S$,\ldots $\bar S_n=(\bar S/S')^{n+1}$,\ldots), and let
$\bar X_\bullet=X'\times_{S'} \bar S_\bullet$. Let
$g_\bullet:\bar X_\bullet \to X'$ be the canonical projection. Put
$E|X'=E'$. The bicomplex ${g_\bullet}_\ast g_\bullet^\ast E'$, whose columns
are the ${g_n}_\ast g_n^\ast E'$, is a resolution of $E'$
\cite[VIII 8]{sga4}. We may therefore represent $C(S',s)$ by a bicomplex whose
columns identify with the $C(\bar S_n,s)$. By (iii),
$C(\bar S_0,s)=C(\bar S,s)$ is acyclic. On the other hand, for $n>0$, we have
$\h^i C(\bar S_n,s)=0$ for $i\leqslant N$ by the induction hypothesis. It
follows that $\h^i C(S',s)=0$ for $i\leqslant N+1$, which completes the proof.





\subsubsection{Remark}\label{VII:5-1-2}

Let $(f,E)$ satisfy condition (i) of \ref{VII:5-1-1}.
\begin{enumerate}[a)]
  \item If $E\in \ob \eD^b(X)$ and $f_\ast$ has finite cohomological
    dimension, then for every $F\in \ob\eD^-(S,A^\circ)$, the canonical arrow
    \begin{equation*}\tag{$*$}\label{VII:eq:5-1-1-1}
      \eR f_\ast(E)\lotimes_A q^\ast F \to \eR f_\ast(E\lotimes_A p^\ast F) \text{,} 
    \end{equation*}
    where $p:X\to S$ and $q:Y\to S$ are the projections, is an isomorphism.
    The same conclusion holds with $E\in\ob\eD^-(X)$.
\end{enumerate} 

Indeed, by dévissage, it is enough to verify the assertion for $F$ of the
form $i_! A_{S'}$, with $i:S'\to S$ quasi-finite. By the Main Theorem, we
reduce to the case where $i$ is finite, and the conclusion is then equivalent
to the commutation of $\eR f_\ast(E)$ with base change by $i$.
\begin{enumerate}[a)]
\setcounter{enumi}{1}
  \item Suppose that $A$ is noetherian and constructible. Then, for every
    constructible $F\in\ob\eD^b(S,A^\circ)$ of finite tor-dimension, the arrow
    \eqref{VII:eq:5-1-1-1} above is an isomorphism.
\end{enumerate}

Indeed, by dévissage, we reduce to the case where $F$ is of the form $i_! M$,
with $i:S'\to S$ locally closed, $A$ locally constant on $S'$, and $M$
constructible, of finite tor-dimension, and with locally constant cohomology
on $S'$. After base change from $S$ to $S'$, we are thus reduced to the case
where $F$ is perfect \cite[I]{sga6}, hence finally, by localization and
dévissage, to the case $F=A$, for which the conclusion is immediate.





\subsubsection{}\label{VII:5-1-3}

Let $f:X\to Y$ be an $S$-morphism and let $E\in\ob\eD^+(X)$. We say that
$(f,E)$ is \emph{cohomologically proper} rel.\ to $S$ if the formation of
$\eR f_\ast E$ commutes with every base change $S'\to S$. Equivalently, after
every base change $S'\to S$, the formation of $\eR f_\ast' E'$ (where
$f'=f\times_S S'$ and $E'=$ the inverse image of $E$ on $X\times_S S'$)
commutes with every finite base change $S'' \to S'$ (cf.
\cite[XII 6.1]{sga4}). Here are some examples of cohomological properness
(for a study of the analogous notion for sheaves of sets or of
noncommutative groups, the reader may consult \cite[XIII]{sga1}).
\begin{enumerate} % originally 1.3.1, 1.3.2, ...
  \item Suppose that $A$ is torsion and $f$ is proper. Then, for every
    $E\in\ob\eD^+(X)$, $(f,E)$ is cohomologically proper rel.\ to $S$
    (proper base-change theorem, \cite[XII 5.1]{sga4}).
  \item Suppose that $S$ is finite, that $X$ and $Y$ are of finite type over
    $S$, and that $A$ is annihilated by an integer invertible on $S$. Then,
    for every $E\in\ob\eD^+(X)$, $(f,E)$ is cohomologically proper rel.\ to
    $S$ (\ref{VII:1-9}).
  \item Suppose that $A$ is constant, noetherian, and annihilated by an
    integer $n$ invertible on $S$, and that
    $f:X=Y\setminus D\to Y$ is the inclusion of the complement of a divisor
    $D$ on $Y$ with normal crossings relative to $S$
    \cite[XIII 2.1]{sga1}. Let $E$ be a sheaf of $A$-modules on $X$ satisfying
    one of the following conditions:
    \begin{enumerate}[(i)]
      \item $E$ is constructible locally constant and tamely ramified along
        $D$;
      \item $E$ is, locally for the étale topology on $Y$ and on $S$, the
        inverse image of a sheaf of $A$-modules on $S$.
    \end{enumerate}
    Then $(f,E)$ is cohomologically proper rel.\ to $S$, and $\eR f_\ast E$
    is constructible (compare \cite[XIII 2.4]{sga1}).
\end{enumerate} 

Let us sketch the proof in case (i). The question is local near a point $y$
of $Y$. We may assume that $D$ is a sum of smooth divisors,
$D=\sum_{i=1}^m D_i$. By the relative Abhyankar lemma
\cite[XIII 5.5]{sga1}, near $y$ we may trivialize $E$ by a finite covering
$\widetilde Y\to Y$ of the form
$\widetilde Y=Y[T_1,\dots,T_r]/(T_1^{n_1}-t_1,\dots,T_r^{n_r}-t_r)$, where the
$t_i$ are local equations of the smooth divisors passing through $y$, and the
$n_i$ are integers prime to the char.\ of $k(y)$. The inverse image
$\widetilde D$ of $D$ in $\widetilde Y$ still has normal crossings rel.\ to $S$, and
if $g:\widetilde X=\widetilde Y\setminus \widetilde D\to X$ is the projection,
$g^\ast E$ is constant. Since $E$ injects into $g_\ast g^\ast E$ and the
quotient is tamely ramified, an easy dévissage reduces us to the case where
$E$ is constant. Denote the projections by $p:X\to S$, $q:Y\to S$, and, for
$1\leqslant i\leqslant m$, let $f_i:Y\setminus D_i \to Y$ be the canonical
inclusion. From the relative purity theorem \cite[XVI 3.7]{sga4}, one easily
deduces by induction on $m$ that, for every locally constant $A$-module $M$ on
$S$, the canonical arrow
\begin{equation*}\tag{5.1.3.1}\label{VII:eq:5-1-3-3-1} % originally 1.3.3.1
  \eR f_\ast (\dZ/n)\lotimes q^\ast M\to \eR f_\ast (p^\ast M) 
\end{equation*}
is an isomorphism, and that on the other hand we have:
\begin{align*}\tag{5.1.3.2}\label{VII:eq:5-1-3-3-2}
  {f_i}_\ast(\dZ/n) &= \dZ/n \text{,} \\
  \eR^1 {f_i}_\ast (\dZ/n) &= (\dZ/n)(-1)_{D_i} \text{,} 
\end{align*}
where a canonical isomorphism is given by the fundamental class $\cl(D_i)$
(\hyperref[IV]{Cycle} \ref{IV:2-1-4}),
\begin{align*}
  \eR^q {f_i}_\ast(\dZ/n) &= 0\text{ for $q>1$,} \\
  \eR f_\ast (\dZ/n) &= \bigotimes_i^\eL \eR {f_i}_\ast (\dZ/n) \text{ , i.e.:} \\
  \eR^1 f_\ast(\dZ/n) &= \bigoplus_{i=1}^m (\dZ/n)(-1)_{D_i} \text{,} \\
  \textstyle\bigwedge^\bullet \eR^1 f_\ast(\dZ/n) &\iso \eR^\bullet f_\ast (\dZ/n) \text{.} 
\end{align*}
The conclusion of 3, in case (i), therefore follows from
\eqref{VII:eq:5-1-3-3-1} and \eqref{VII:eq:5-1-3-3-2}. In case (ii), by
passage to the limit we reduce to the case where $E=p^\ast M$, with $M$
constructible, and then, by dévissage and base change (using (i)), to the
already treated case where $M$ is constant.

It follows from the preceding proof that \eqref{VII:eq:5-1-3-3-1} is an
isomorphism for every $M\in\ob\eD^+(S)$.

One may compare the formulas \eqref{VII:eq:5-1-3-3-2} with the analogous
formulas in integral cohomology when $S=\spec(\dC)$ (see, for example,
\cite[3.1]{de71-2}), which follow from the fact that, locally for the classical
topology, the complement of $D$ has the homotopy type of a torus.

By passage to the limit, one deduces from \eqref{VII:eq:5-1-3-3-2} similar
formulas in $\ell$-adic cohomology, i.e. with $\dZ/n$ replaced by $\dZ_\ell$
($\ell$ a prime invertible on $S$). Suppose that $q:Y\to S$ is proper and
smooth. It then follows from the Weil conjectures \cite{de74,de80} that the
Leray spectral sequence
\[
  E_2^{i j} = \eR^i q_\ast \eR^j f_\ast \dZ_\ell \Rightarrow \eR^\bullet p_\ast \dZ_\ell 
\]
degenerates at $E_3$ modulo torsion \cite[\S 6]{de71-1}. The same phenomenon
occurs in integral cohomology when $S=\spec(\dC)$; see (loc.\ cit.) and
\cite[3.2.13]{de71-2}.





\subsection{Specialization and cospecialization}\label{VII:5-2}

Let $f:X\to S$ and $E\in\ob\eD^+(X)$.






\subsubsection{}\label{VII:5-2-1}

Suppose that the formation of $\eR f_\ast E$ commutes with every finite base
change. For every specialization $u:t\to S(s)$ of geometric points of $S$,
we then define an arrow, called the \emph{specialization arrow},
\begin{equation*}\tag{5.2.1.1}\label{VII:eq:5-2-1-1}
  \spe(u):\eR\Gamma(X_s,E|X_s) \to \eR\Gamma(X_t,E|X_t) \text{,}
\end{equation*}
as follows. Let, as in \ref{VII:5-1-1}(iii), $\bar S$ be the normalization in
$k(t)$ of the closure of $u(t)$ in $S(s)$, and let
$\bar f:\bar X\to \bar S$ be the morphism deduced from $f$ by the base change
$\bar S\to S$. By the immediate implication (i)$\Rightarrow$(iii) of
\ref{VII:5-1-1}, the base-change arrow
\begin{equation*}\tag{5.2.1.2}\label{VII:eq:5-2-1-2}
  \eR\Gamma(\bar X,E|\bar X)=\eR\bar f_\ast (E|\bar X)_s \to \eR\Gamma(X_s,E|X_s) 
\end{equation*}
is an isomorphism. We define \eqref{VII:eq:5-2-1-1} as the composite of the
inverse of \eqref{VII:eq:5-2-1-2} with the restriction arrow
\begin{equation*}\tag{5.2.1.3}\label{VII:eq:5-2-1-3}
  \eR\Gamma(\bar X,E|\bar X) \to \eR\Gamma(X_t,E|X_t) \text{.}
\end{equation*}






\subsubsection{}\label{VII:5-2-2}

Suppose that $f$ is locally acyclic rel.\ to $E$ (\ref{VII:2-12}). For
every specialization $u:t\to S(s)$ of geometric points of $S$, we then define
an arrow, called the \emph{cospecialization arrow},
\begin{equation*}\tag{5.2.2.1}\label{VII:eq:5-2-2-1}
  \cosp(u):\eR\Gamma(X_t,E|X_t) \to \eR\Gamma(X_s,E|X_s) \text{,} 
\end{equation*}
as follows. Let $g:X_t\to X$ and $g':\bar X\to X$ be the canonical arrows.
Viewed vis-\`a-vis these two maps, the canonical arrow
\begin{equation*}\tag{5.2.2.2}\label{VII:eq:5-2-2-2}
  \eR g'_\ast(E|\bar X)\to \eR g_\ast(E|X_t) 
\end{equation*}
is an isomorphism by the local-acyclicity hypothesis. Indeed, we may assume
that $S$ is strictly local and integral, with closed point $s$, and that $t$
is a geometric generic point. Compute the fiber of
\eqref{VII:eq:5-2-2-2} at a geometric point $x$ of $X$, with image $y$ in
$S$. After replacing $S$ by the strict localization of $S$ at $y$, we may
assume that $y=s$. Since $\bar S\to S$ is an inverse limit of finite
morphisms and is radicial at $s$, we find that
$\eR g'_\ast(E|\bar X)_x = E_x$ and that the fiber of
\eqref{VII:eq:5-2-2-2} at $x$ identifies with the restriction
$E_x=\eR\Gamma(X(x),E)\to \eR\Gamma(X(x)_t,E|X(x)_t)$, which is an
isomorphism by hypothesis. Applying $\eR\Gamma(X,-)$ to
\eqref{VII:eq:5-2-2-2}, we deduce that the restriction arrow
\eqref{VII:eq:5-2-1-3} is an isomorphism. We define
\eqref{VII:eq:5-2-2-1} as the composite of the inverse of
\eqref{VII:eq:5-2-1-3} with the base-change arrow
\eqref{VII:eq:5-2-1-2}.






\subsubsection{}\label{VII:5-2-3}

If $s'' \xrightarrow{u'} s' \xrightarrow u s$ are specialization arrows of
geometric points of $S$, one verifies that
\begin{align*}
  \spe(u')\spe(u) &= \spe(u u') \text{,} \\
  \cosp(u) \cosp(u') &= \cosp(u u') \text{,} 
\end{align*}
whenever both sides are defined.






\subsubsection{}\label{VII:5-2-4}

Suppose that the formation of $\eR f_\ast E$ commutes with every finite base
change and that $f$ is locally acyclic rel.\ to $E$. It follows from the
definitions that, for every specialization $u:t\to S(s)$ of geometric points
of $S$, the arrows $\spe(u)$ and $\cosp(u)$ are isomorphisms, inverse to one
another.

The preceding hypotheses hold in particular when $A$ is noetherian and
annihilated by an integer $n$ invertible on $S$, when $f$ is proper and smooth,
and when $E$ has locally constant cohomology (cohomological properness of
proper morphisms (\ref{VII:5-1-3}), and local acyclicity of smooth morphisms
\cite[XV]{sga4}). This recovers the ``specialization theorem''
\cite[XVI 2.1]{sga4}.

On the other hand, \ref{VII:5-1-1} has the following consequence:






\begin{proposition}\label{VII:5-2-5}
Suppose that $f$ is locally acyclic relative to $E$ and that, for every
specialization $u:t\to S(s)$, the cospecialization arrow $\cosp(u)$ is an
isomorphism. Then the formation of $\eR f_\ast E$ commutes with every finite
base change.
\end{proposition}

In view of \ref{VII:5-2-4}, \ref{VII:5-2-5} generalizes
(\hyperref[I]{Étale cohomology} \ref{I:5-1-7}).






\begin{corollary}\label{VII:5-2-6}
Suppose that $f$ is locally acyclic rel.\ to $E$. Let $x$ be a geometric
point of $X$, with image $s$ in $S$, and denote by
$f_{(x)}:X(x)\to S(s)$ the morphism induced by $f$. Then the formation of
$\eR {f_{(x)}}_\ast(E|X(x))$ commutes with every finite base change
$S'\to S(s)$.
\end{corollary}

By \ref{VII:5-2-5}, it is enough to verify that the cospecialization arrows
for $(f_{(x)},E|X(x))$ are isomorphisms. By the transitivity of the
cospecialization arrows (\ref{VII:5-2-3}), it is enough to show that if
$u:t\to S(s)$ is a specialization, then
$\cosp(u):\eR\Gamma(X(x)_t,E|X(x)_t)\to \eR\Gamma(X(x)_s,E|X(x)_s)$ is an
isomorphism. But since $s$ is closed in $S(s)$, $X(x)_s$ is strictly local
with closed point $x$, $\eR\Gamma(X(x)_s,E|X(x)_s)=E_x$, and $\cosp(u)$ is an
isomorphism, inverse to the restriction isomorphism
$E_x\to \eR\Gamma(X(x)_t,E|X(x)_t)$.






\begin{corollary}\label{VII:5-2-7}
Let $f:X\to Y$ be an $S$-morphism, let $p:X\to S$, $q:Y\to S$ be the
projections, and let $E\in\ob\eD^+(X)$. Suppose that $A$ is locally constant,
that $f$ is locally acyclic rel.\ to $E$, and that $q$ is locally acyclic at
every geometric point $y$ of $Y$ rel.\ to every sheaf locally constant
in a neighborhood of $y$. Then $p$ is locally acyclic relative to $E$.
\end{corollary}

Let $x$ be a geometric point of $X$, with images $y$ and $s$ in $Y$ and $S$,
let $t\to S(s)$ be a specialization, and let
$f_{(x)}:X(x)\to Y(y)$ be the morphism induced by $f$. We must show that the
restriction
\begin{equation*}\tag{$*$}\label{VII:eq:5-2-7-1}
  E_x \to \eR\Gamma(X(x)_t,E|X(x)_t) 
\end{equation*}
is an isomorphism. We may assume that $A$ is constant. By \ref{VII:5-2-4} and
\ref{VII:5-2-5}, for every specialization $z\to Y(y)$, the specialization
arrow
\[
  E_x = \eR {f_{(x)}}_\ast (E|X(x))_y \to (\eR {f_{(x)}}_\ast (E|X(x)))_z = \eR\Gamma(X(x)_z,E|X(x)_z) 
\]
is an isomorphism; in other words, $\eR{f_{(x)}}_\ast(E|X(x))$ is constant
with value $E_x$. The acyclicity hypothesis on $q$ therefore plays the
decisive r\^ole and ensures that the restriction
\[
  E_x \to \eR\Gamma(Y(y)_t,(E_x)|Y(y)_t) = \eR\Gamma\left(Y(y)_t,\eR {f_{(x)}}_\ast (E|X(x))|Y(y)_t\right)
\]
is an isomorphism. But this arrow identifies with
\eqref{VII:eq:5-2-7-1}, which proves the assertion.

The hypothesis of \ref{VII:5-2-7} on $q$ holds in particular when $q$ is
smooth and $A$ is annihilated by an integer $n$ invertible on $S$
\cite[XV]{sga4}. On the other hand, if the hypotheses of \ref{VII:5-2-7} hold
after every base change $S'\to S$, then it follows that $p$ is universally
locally acyclic rel.\ to $E$. These remarks imply, in particular, the
statement of \ref{VII:2-14}.






\begin{corollary}\label{VII:5-2-8}
Let $f:X\to S$ and $E\in \ob\eD^+(X)$. Suppose that $f$ is universally locally
acyclic rel.\ to $E$, that the formation of $\eR f_\ast E$ commutes with
every finite base change, and that, for every morphism $s'\to s$, where $s$ is
a geometric point of $S$ and $s'$ is the spectrum of a separably closed
field, the base-change arrow
$\eR\Gamma(X_s,E|X_s)\to \eR\Gamma(X_{s'},E|X_{s'})$ is an isomorphism. Then
$(f,E)$ is cohomologically proper (\ref{VII:5-1-3}).
\end{corollary}

Let $f':X'\to S'$ be the morphism deduced from $f$ by a base change
$S'\to S$; we must show that the formation of $\eR f'_\ast(E|X')$ commutes
with every finite base change. By \ref{VII:5-2-5}, it is enough to show that,
for every specialization $u':t'\to S'(s')$, the cospecialization arrow
$\cosp(u')$:
$\eR\Gamma(X_{t'}',E|X_{t'}') \to \eR\Gamma(X_{s'}',E|X_{s'}')$ is an
isomorphism. Let $s$ (resp.\ $t$) be the geometric point that is the image of
$s'$ (resp.\ $t'$) in $S$, and let $u:t\to S(s)$ be the
specialization defined by $u'$. One readily sees that there is a commutative
square
\[\xymatrix{
  \eR\Gamma(X_t,E|X_t) \ar[r]^-{\cosp(u)} \ar[d] 
    & \eR\Gamma(X_s,E|X_s) \ar[d] \\
  \eR\Gamma(X_{t'}',E|X_{t'}') \ar[r]^-{\cosp(u')} 
    & \eR\Gamma(X_{s'}',E|X_{s'}') 
}\]
in which the vertical arrows are the base-change arrows. By hypothesis, these
are isomorphisms. On the other hand, (\ref{VII:5-2-4}) $\cosp(u)$ is an
isomorphism, hence $\cosp(u')$ is an isomorphism, which completes the proof.

Here is an example of an application of \ref{VII:5-2-8}. Suppose that $A$ is
annihilated by an integer $n$ invertible on $S$, and that $f$ is locally of
finite type and universally locally acyclic. Let $x$ be a geometric point of
$X$, with image $s$ in $S$, and let $f_{(x)}:X(x)\to S(s)$ be the morphism
induced by $f$. It follows from \ref{VII:5-2-6} and \ref{VII:1-9} (by passage
to the limit) that the hypotheses of \ref{VII:5-2-8} are satisfied by the
pair $(f_{(x)},E|X(x))$, which is therefore cohomologically proper.






\subsubsection{}\label{VII:5-2-9}

Suppose that $A$ is constant. Let $f:X\to S$, and let $E\in\ob\eD^+(X)$ be a
complex of finite tor-dimension. We say that $f$ is \emph{strongly locally
acyclic} rel.\ to $E$ if, for every geometric point $x$ of $X$, with image
$s$ in $S$, every specialization $t\to S(s)$, and every $A^\circ$-module $M$,
the restriction arrow
\[
  E_x\lotimes_A M \to \eR\Gamma(X(x)_t,(E|X(x)_t)\lotimes_A M) 
\]
is an isomorphism. I do not know whether ``locally acyclic'' implies
``strongly locally acyclic.''






\begin{proposition}\label{VII:5-2-10}
Under the hypotheses of \ref{VII:5-2-9}, suppose that $A$ is noetherian and
torsion, that $S$ is noetherian, and that $f$ is strongly locally acyclic
rel.\ to $E$. Then, for every diagram with cartesian squares
\[\xymatrix{
  X \ar[d]^-f 
    & \ar[l]_-j X' \ar[d]^-{f'} 
    & \ar[l]_-{j'} X'' \ar[d]^-{f''} \\
  S 
    & \ar[l]_-i S' 
    & \ar[l]_-{i'} S'' \text{,} 
}\]
with $i$ quasi-finite and $i'$ an open immersion, and every
$F\in\ob\eD^+(S'',A^\circ)$, the canonical arrow
\begin{equation*}\tag{$*$}\label{VII:eq:5-2-10-1}
  j^\ast E \lotimes {f'}^\ast \eR i'_\ast F \to \eR j'_\ast ((j j')^\ast E\lotimes {f''}^\ast F) 
\end{equation*}
is an isomorphism.
\end{proposition}

Compare this statement with \cite[XV 1.17]{sga4}.

Let us prove \ref{VII:5-2-10}. We readily reduce to the case where $S=S'$ is
strictly local, and to $F$ constructible and concentrated in degree $0$. We then
resolve $F$ on the right by finite sums of sheaves of the form $u_\ast M$,
where $u:t\to S$ is a geometric point and $M$ is an $A^\circ$-module. The
hypothesis, vis-\`a-vis this resolution, implies that
\eqref{VII:eq:5-2-10-1} is an isomorphism for $F=u_\ast M$, and the result
follows.
